\usepackage[utf8]{inputenc} % Pacote para permitir o uso de acentos
% Ajustes de hifenização
\tolerance=1
\emergencystretch=\maxdimen
\hyphenpenalty=10000
\hbadness=10000
\usepackage{lipsum} % Pacote para gerar texto aleatório
\usepackage{ragged2e} % Pacote para mais opções de alinhamento
\usepackage{graphicx} % Pacote para inserir figuras jpeg, pdf e png
\usepackage{ae}
\usepackage{epsfig}
\usepackage{mathtools}
\usepackage[export]{adjustbox}
\usepackage{lineno}
\usepackage{amsthm}
\usepackage{commath}
\usepackage{color, colortbl}
\usepackage{bbm}
\usepackage{setspace} % Pacote para ajustar espaçamento entre linhas

% Instructions to prevent LaTeX from hyphenating words
\tolerance=1
\emergencystretch=\maxdimen
\hyphenpenalty=10000
\hbadness=10000

\graphicspath{{figures/}} % Indica o diretório no qual as imagens estão sendo salvas, evitando que precisamos digitar parte do caminho da pasta que contém as imagens
\usepackage[rightcaption]{sidecap} % Pacote para mais opções de figura, colocar, título da legenda da figura ao lado direito
\usepackage{url} % Pacote para citar referências de sites
\usepackage{makeidx} % Pacote para criar índice
\makeindex % Habilitar a criação do índice
\usepackage{longtable} % Pacote para usar tabelas longas

% ----------
% Pacotes para dar quebra de linha dentro da tabela:
%\usepackage{fourier} 
%\usepackage{array}
\usepackage{makecell}

\renewcommand\theadalign{bc}
\renewcommand\theadfont{\bfseries}
\renewcommand\theadgape{\Gape[4pt]}
\renewcommand\cellgape{\Gape[4pt]}
% ----------

\usepackage{wrapfig} % Pacote para alinhar figuras à esquerda ou à direita do texto
\usepackage{ae}
\usepackage{epsfig}
\usepackage{mathtools}
\usepackage[export]{adjustbox}
\usepackage{lineno}
\usepackage{amsfonts}
\usepackage{amsthm}
\usepackage{commath}
\usepackage{color, colortbl}
\usepackage{bbm}

% ----------
% Pacote para inserir cabeçalho e rodapé:
% Atenção! Se remover o cabeçalho e rodapé, o LaTeX coloca como cabeçalho o padrão dele, que é bem interessante também!
% \usepackage{fancyhdr} % Pacote para inserir cabeçalho e rodapé
% \pagestyle{fancy} % Estilo do cabeçalho e rodapé, neste caso apenas uma linha no cabeçalho
%\lhead{{\footnotesize Prandiano\textregistered}} % Inserir texto no canto esquerdo do cabeçalho
%\chead{centro} % Inserir texto no centro do cabeçalho
%\rhead{{\footnotesize Curso de Python}} % Inserir texto no canto direito do cabeçalho
%\lfoot{{\footnotesize Copyright \textcopyright 2018 Prandiano}} % Inserir texto no canto esquerdo do rodapé
%\cfoot{centro} % Inserir texto no centro do rodapé
%\rfoot{{\footnotesize Reprodução proibida: Lei 9.610 de 19/02/1998}} % Inserir texto no canto direito do rodapé
% O comando \renewcommand re-nova o parâmetro especificado
%\renewcommand{\footrulewidth}{0.4pt} % Inserir uma linha acima do rodapé
% ----------

% ----------
% Pacote para inserir múltiplas figuras:
\usepackage{subfigure} % Pacote para inserir múltiplas figuras
% \usepackage{subcaption} % Pacote incompatível com subfigure
%\usepackage{caption}
% ----------

\usepackage{multicol} % Pacote que permite usar multicolunas
\usepackage{float} % Pacote para usar o "H" do figure
\usepackage[dvipsnames]{xcolor} % Pacote de cores para textos e possivelmente outros
\usepackage{amsmath} % Pacote para usar o cases
\usepackage[makeroom]{cancel} % Pacote para usar o cancelar (em e embaixo)
\usepackage{enumitem} % Pacote para enumerar itens
\usepackage{amssymb} % Pacote para inserir os símbolos dos conjuntos numéricos N, Z, Q, R, C
\usepackage[T1]{fontenc}
\usepackage{enumitem}
\usepackage{textcomp}
\usepackage{pdfpages} % Pacote para inserir página em PDF
\usepackage{courier} % Pacote para usar a fonte courier
%\usepackage{matlab-prettifier}
\usepackage{lscape} % Pacote para mudar sentido da página
\newtheorem{theorem}{Teorema} % Pacote para inserir Teorema
\newtheorem{corollary}{Corollary} % Pacote para inserir Corolário
\newtheorem{remark}{Remark}
\newtheorem{mytheorem}{Definition}
\usepackage{hyperref} % Pacote para inserir hyperref
\usepackage{multirow} % Pacote para usar os comandos \multirow
\usepackage{pdfpages} % Pacote para incluir PDF's
\usepackage{listingsutf8} % Pacote para inserir arquivos externos tipo .txt
\usepackage{etoolbox}
\AtBeginEnvironment{verbatim}{\footnotesize}

\usepackage{standalone} % Pacote para inserir sub-arquivos LaTeX
\usepackage{bookmark} % Pacote para NÃO tratar o Índice como Apêndice e colocar letra

% Nome da listagem em português
\renewcommand{\lstlistingname}{Listagem}

% REFERÊNCIAS: ----------
%\usepackage[numbers, compress]{natbib} % Pacote para usar as referências bibliográficas
%\usepackage{usebib}
% \bibliographystyle{references} % Estilo das referências
%\setcitestyle{authoryear,aysep={},open={(},close={)}}
%\newcommand{\myand}{e}

% ----------


% PARÂMETROS DA FUNÇÃO LISTINGS: ----------
\lstset
{
    basicstyle=\small\ttfamily,
    commentstyle=\color{Green},
    keywordstyle=\color{Cerulean},
    frame=single,
    language=Python,
    morekeywords={True, False},
    numbers=left,
    numbersep=10pt,
    numberstyle=\footnotesize\color{Gray},
    showstringspaces=false,
    stringstyle=\color{Mulberry},
    tabsize=3,
    breaklines=true, % Este comando é para quebrar uma linha comprida
    postbreak=\mbox{\textcolor{red}{$\hookrightarrow$}\space}, % Este comando é para indicar com uma seta vermelha que a linha continua embaixo 
    extendedchars=true,
    literate=
        {á}{{\'a}}1 {Á}{{\'A}}1 {à}{{\`a}}1 {À}{{\`A}}1 {ã}{{\~a}}1 {Ã}{{\~A}}1 {â}{{\^a}}1 {Â}{{\^A}}1
        {ç}{{\c{c}}}1 {Ç}{{\c{C}}}1 {é}{{\'e}}1 {É}{{\'E}}1 {ê}{{\^e}}1 {Ê}{{\^E}}1 {í}{{\'\i}}1 {Í}{{\'I}}1
        {ó}{{\'o}}1 {Ó}{{\'O}}1 {ô}{{\^o}}1 {Ô}{{\^O}}1 {õ}{{\~o}}1 {Õ}{{\~O}}1 {ú}{{\'u}}1 {Ú}{{\'U}}1
        {ü}{{\"u}}1 {Ü}{{\"U}}1 {“}{{``}}1 {”}{{''}}1 {’}{{'}}1 {—}{{--}}1 {–}{{--}}1 {−}{{-}}1
        {±}{{$\pm$}}1 {σ}{{$\sigma$}}1 {→}{{$\to$}}1 {…}{{\ldots}}1
}

\lstset{inputencoding=utf8/latin1}

\lstdefinelanguage{Matlab}{
	language=Matlab,
	basicstyle=\ttfamily\small,
	keywordstyle=\color{blue}\bfseries,
	stringstyle=\color{red},
	commentstyle=\color{gray}\itshape,
	morekeywords={function, end, if, else, for, while, switch, case, otherwise},
	sensitive=true
}
% ----------

% ----------
\usepackage{nomencl} % Pacote para criar nomenclatura e lista de nomenclatura
\makenomenclature % Criar a nomenclatura
\renewcommand{\nomname}{Lista de Nomenclaturas} % Renomear o título da nomenclatura
% ----------

\usepackage{amsmath} % Pacote para inserir matrix com chaves, entre outras coisas

\usepackage{booktabs} % Pacote para usar \toprule, \midrule e \bottomrule


% EXECUTAR CÁLCULOS A PARTIR DA LINHA DE COMANDO: ---------
%\usepackage{siunitx, xfp} % Pacote para executar cálculos a partir da linha de comando
%\usepackage{l3regex}%unless you have very new expl3
%%\sisetup{locale=DE}
%%\let\oldnum\num
%%\ExplSyntaxOn
%\def\num#1{\def\tmp{#1}\regex_replace_all:nnN{,([0-9])}{.\1}{\tmp}\oldnum{\tmp}}
%\ExplSyntaxOff
% ----------


%\pagecolor{black} % Alterar a cor da página
%\color{white} % Alterar a cor do texto

% ----------
% Pacotes para inserir algoritmos e pseudo-códigos:
\setlength{\marginparwidth}{2.5cm}
\usepackage[colorinlistoftodos]{todonotes}
\usepackage{algorithm} % Pacote para inserir algoritmos
\usepackage{algpseudocode} % Pacote para inserir algoritmos e pseudo-códigos
\floatname{algorithm}{Algoritmo} % Renomear a palavra Algorithm para Procedimento
\renewcommand{\algorithmicrequire}{\textbf{entrada:}} % Renomear a palavra Input para Entrada
\renewcommand{\algorithmicensure}{\textbf{saída:}} % Renomear a palavra Output para Saída
\renewcommand{\algorithmicwhile}{\textbf{enquanto}} % Renomear a palavra While para Enquanto
\renewcommand{\algorithmicif}{\textbf{se}} % Renomear a palavra If para Se
\renewcommand{\algorithmicend}{\textbf{fim}} % Renomear a palavra End Output para Fim do
\renewcommand{\algorithmicprocedure}{\textbf{procedimento}} % Renomear a palavra Procedure para Procedimento
\renewcommand{\algorithmicfor}{\textbf{para}} % Renomear a palavra For para Para
\renewcommand{\algorithmicelse}{\textbf{senão}} % Renomear a palavra Else para Senão
\renewcommand{\algorithmicfunction}{\textbf{função}} % Renomear a palavra Function para Função 
\renewcommand{\algorithmicthen}{\textbf{então}} % Renomear a palavra Then para Então
\renewcommand{\algorithmicdo}{\textbf{fazer}} % Renomear a palavra Do para Fazer
% ----------

\newcommand{\quotes}[1]{``#1''} % Comando para inserir um texto dentro de aspas-dupla com espaço após o final

% FLUXOGRAMA: ----------
\usepackage{tikz}
\usetikzlibrary{positioning, shapes.geometric}
\tikzstyle{line}=[draw] % here

\tikzstyle{startstop} = [rectangle, rounded corners, minimum width=3cm,
minimum height=1cm,text centered, draw=black, fill=white!30]
\tikzstyle{arrow} = [thick, ->, >=stealth]
\tikzstyle{process} = [rectangle, minimum width=3cm, minimum height=1cm, text centered, draw=black, fill=white!30]
\tikzstyle{io} = [trapezium, trapezium left angle=70, trapezium right angle=110, minimum width=3cm, minimum height=1cm, text centered, draw=black, fill=white!30]
\tikzstyle{connector} = [draw, circle, fill=white!20]
\tikzstyle{decision} = [diamond, minimum width=3cm, minimum height=1cm, text centered, draw=black, fill=white!30, aspect=3]
% ----------


% CAIXA ENTORNO DAS FÓRMULAS: ----------
\usepackage{dashbox}
\newcommand\dboxed[1]{\dbox{\ensuremath{#1}}}

% ----------

%\usepackage{tabularx} % para colunas com largura ajustável

% -------- CONFIGURAÇÃO DE MARGENS --------
% \usepackage[top=2cm,bottom=2cm,left=3cm,right=2cm]{geometry}

% -------- PACOTE fancyhdr --------
\usepackage{fancyhdr}

\setlength{\marginparwidth}{2cm}

% Ajuste do tamanho do cabeçalho para acompanhar exatamente \textwidth
\setlength{\headwidth}{\textwidth}

% Define um headheight mínimo (para evitar avisos)
\setlength{\headheight}{15pt} 

% Limpa configurações padrão de cabeçalho/rodapé
\pagestyle{fancy}
\fancyhf{} 

% Marcações automáticas de seção para o cabeçalho (opcional)
\renewcommand{\sectionmark}[1]{%
	\markboth{\thesection\ #1}{}%
}

% Texto do cabeçalho à esquerda: seção e título
% Usamos \leftmark pois o \sectionmark acima grava no \leftmark
\fancyhead[L]{\small \leftmark}

% Texto do cabeçalho à direita: número de página
\fancyhead[R]{\small \thepage}

% Linha horizontal no cabeçalho (espessura 0.4pt)
\renewcommand{\headrulewidth}{0.4pt}

% Sobrescreve o título padrão “References” para “Referências”
%\renewcommand{\refname}{Referências}
% Tradução de autoref para o português
\def\figureautorefname{Figura}
\def\tableautorefname{Tabela}
\def\equationautorefname{Equa\c{c}\~ao}

% \hypersetup{backref=true}

\ifdefined\backref
        \renewcommand*{\backref}[1]{}
        \renewcommand*{\backrefentrycount}[2]{#1\ifnum#2>1 ~(#2)\fi}
\fi

\DeclarePairedDelimiter{\ceil}{\lceil}{\rceil}

\onehalfspacing

\usepackage{pgfgantt}
